\حصہ{مسئلہ  پھیلاؤ اور ایک  بنیادی مسئلہ}
مستوی میں  مسئلہ گرین  کہتا  ہے کہ   سادہ بند منحنی کے پار  سمتی میدان کا  خالص بیرونی بہاؤ     منحنی میں محیط خطے پر میدان کے پھیلاؤ  پر تکمل  سے حاصل کیا جا سکتا ہے۔ تین ابعاد  میں متعلقہ  مسئلہ، جو مسئلہ پھیلاؤ کہلاتا ہے،  کے تحت فضا میں بند سطح سے   سمتیہ میدان  کا  باہر رخ گزرتا بہاؤ سطح میں محیط  خطے پر  میدان کے پھیلاؤ    کے تکمل سے حاصل  کیا جا سکتا ہے۔اس حصے میں ہم مسئلہ پھیلاؤ ثابت کرتے ہیں، اور دکھاتے ہیں یہ بہاؤ کی قیمت کی  تلاش  کیسے آسان بناتا ہے۔ ہم  برقی میدان میں بہاؤ  کا قانون گاوس اور  ماقوا حرکیات  کی  استمراری مساوات بھی اخذ کرتے ہیں۔

\جزوحصہء{تین ابعاد میں پھیلاؤ}
سمتی میدان \عددی{\kvec{F} = M(x,y,z)\ai + N(x,y,z)\aj + P(x,y,z)\ak}  کا\اصطلاح{ پھیلاؤ}\فرہنگ{پھیلاؤ}\حاشیہب{divergence}\فرہنگ{divergence} درج ذیل غیر سمتی تفاعل ہو گا۔
\begin{align}
\kvec{F}\text{پھیلاؤ}  = \nabla \cdot \kvec{F}=\frac{\partial M}{\partial x}+\frac{\partial N}{\partial y}+\frac{\partial P}{\partial z}\tag{1}
\end{align}
علامت \قول{پھیلاؤ $\kvec{F}$} کو  \قول{$F$ کا  پھیلاؤ} پڑھا جاتا ہے، جبکہ علامت  \قول{$\nabla\cdot \kvec{F}$}  کو  \قول{ڈیل  ایف} پڑھا جاتا ہے۔

 تین ابعاد میں  پھیلاؤ  $\kvec{F}$ کا  وہی طبیعی   مفہوم  ہے جو دو ابعاد  میں ہے۔ اگر  \عددی{\kvec{F}} سیال کے بہاؤ کا سمتی میدان رفتار  ہو،  نقطہ  \عددی{(x,y,z)}  پر  بہاؤ \عددی{\kvec{F}}کی قیمت   اس نقطے پر سیال  کے دخول یا خروج کی شرح  دے گا۔  پھیلاؤ سے مراد  نقطے پر بہاؤ  فی اکائی حجم   یعنی کثافت بہاؤ ہے۔


\ابتدا{مثال}\موٹا{پھیلاؤ کی تلاش}

سمتی میدان \عددی{\kvec{F} = 2xz\ai - xy\aj - z\ak} کا پھیلاو تلاش کریں۔

حل: \quad
سمتی میدان \عددی{\kvec{F}} کا  پھیلاؤ درج ذیل ہے۔
\begin{align*}
\nabla \cdot \kvec{F}=\frac{\partial}{\partial x}(2xz)+\frac{\partial}{\partial y}(-xy)+\frac{\partial}{\partial z}(-z)=2z-x-1
\end{align*}
\انتہا{مثال}


\جزوحصہء{مسئلہ پھیلاؤ}

مسئلہ پھیلاؤ  کہتا ہے مناسب شرائط کے تحت،       (باہر  رخ سمت بند)  بند سطح سے  سمتی میدان کا  باہر  گزرتا بہاؤ سطح میں محیط خطہ پر میدان کے پھیلاؤ کے تہرا تکمل کے  برابر ہو گا۔


\ابتدا{مسئلہ}\موٹا{مسئلہ پھیلاؤ}

 بند ، سمت بند سطح \عددی{S}    سے   سمتی میدان \عددی{\kvec{F}} کا،      سطح کے باہر رخ  اکائی عمودی سمتیہ کے رخ ،گزرتا بہاؤ،  سطح میں محیط خطہ \عددی{D}  پر \عددی{\nabla\cdot \kvec{F}} کے تکمل کے برابر ہو گا۔
\begin{align}
\underbrace{\iint_S \kvec{F}\cdot\kvec{n}\dif \sigma}_{\text{\RL{باہر رخ بہاؤ}}}=\underbrace{\iint_D \nabla \cdot\kvec{F}\dif H}_{\text{\RL{تکمل پھیلاؤ}}}\tag{2}
\end{align}
\انتہا{مسئلہ}

\ابتدا{مثال}
کرہ \عددی{x^2+y^2+z^2=a^2} پر سمتی میدان \عددی{\kvec{F}=x\ai+y\aj+z\ak}  کے لئے مساوات \حوالہء{2} کے دونوں اطراف حل کریں۔

حل:\quad
سطح \عددی{S} کا ( درج ذیل)  باہر رخ اکائی عمودی سمتیہ  \عددی{f(x,y,z)=x^2+y^2+z^2-a^2} کے ڈھلوان  سے حاصل کرتے ہیں۔	
\begin{align*}
\an=\frac{2(x\ai+y\aj+z\ak)}{\sqrt{4(x^2+y^2+z^2)}}=\frac{x\ai+y\aj+z\ak}{a}
\end{align*}
یوں  درج ذیل ہو جہاں سطح پر \عددی{x^2+y^2+z^2=a^2} لیا گیا ہے۔
\begin{align*}
\kvec{F}\cdot \an \dif\sigma=\frac{x^2+y^2+z^2}{a}\dif \sigma=\frac{a^2}{a}\dif\sigma=a\dif\sigma
\end{align*}
لہٰذا مساوات \حوالہء{2} کا ایک ہاتھ   ا درج ذیل ہو گا۔
\begin{align*}
\iint_S \kvec{F}\cdot \an\dif \sigma=\iint_S a\dif \sigma=a\iint_S \dif\sigma=4\pi a^3
\end{align*}
میدان \عددی{\kvec{F}} کا پھیلاؤ 
\begin{align*}
\nabla\cdot \kvec{F}=\frac{\partial}{\partial x}(x)+\frac{\partial}{\partial y}(y)+\frac{\partial}{\partial z}(z)=3
\end{align*}
لہٰذا  مساوات \حوالہء{2} کا دوسرا ہاتھ ا درج ذیل ہو گا۔
\begin{align*}
\iiint_D \nabla\cdot \kvec{F}\,\dif H=\iiint_D 3 \dif H=3\left(\frac{4}{3}\pi a^3\right)=4\pi a^3
\end{align*}
\انتہا{مثال}

\ابتدا{مثال}\موٹا{بہاؤ کا حصول}

ثُمن اول سے مستویات \عددی{x=1}، \عددی{y=1}، اور \عددی{z=1}  مکعب کاٹتی ہیں۔  سمتی میدان \عددی{\kvec{F}=xy\ai+yz\aj+xz\ak} کا  مکعب کی سطح سے باہر رخ گزرتا بہاؤ تلاش کریں۔

حل:\quad
مکعب کے چھ  سطوح پر تکمل لے کر بہاؤ  کی قیمت تلاش کرنے کی بجائے ہم  پھیلاؤ
\begin{align*}
\nabla\cdot\kvec{F}=\frac{\partial}{\partial x}(xy)+\frac{\partial}{\partial y}(yz)+\frac{\partial}{\partial z}(xz)=y+z+x
\end{align*}
 کا تکمل   مکعب کے اندرون پر لیتے ہیں۔
\begin{align*}
\text{بہاؤ}&=\iint_{\text{\RL{کعبی سطوح}}}\kvec{F}\cdot\an\dif\sigma=\iiint_{\text{\RL{کعبی اندرون}}}\nabla \cdot\kvec{F}\dif H&&\text{\RL{مسئلہ پھیلاو}}\\
&=\int_0^1\int_0^1\int_0^1 (x+y+z)\dif x\dif y\dif z=\frac{3}{2}
\end{align*}
\انتہا{مثال}

%KKK??? am here

\جزوحصہء{مخصوص خطوں  کے لئے مسئلہ پھیلاؤ کا ثبوت}
ہم فرض کرتے ہیں کہ \عددی{\kvec{F}} کے اجزاء کے جزوی یک رتبی  تفرق استمراری ہیں۔ ہم مزید فرض کرتے ہیں کہ خطہ  \عددی{D} محدب ہے اور اس میں کوئی سوراخ  یا بلبلا نہیں پایا جاتا، مثلاً ٹھوس مکعب، کرہ، یا ترخیمی جسم، اور \عددی{S} ٹکڑوں میں ہموار سطح  ہے۔ اس کے ساتھ ساتھ ہم فرض کرتے ہیں کہ \عددی{xy} مستوی پر \عددی{D} کا تظلیل قائمہ \عددی{R_{xy}} ہے   اور اس کے  اندر \عددی{xy}  مستوی کو ہر عمودی لکیر  سطح \عددی{S} کو ٹھیک دو نقاط  پر قطع کرتی ہے جس سے ہمیں درج ذیل دو سطوح ملتے ہیں
\begin{align*}
S_1:\quad z&=f_1(x,y),\quad \text{\RL{$(x,y)$ تظلیل $R_{xy}$ میں ہیں}}\\
S_2:\quad z&=f_2(x,y),\quad \text{\RL{$(x,y)$ تظلیل $R_{xy}$ میں ہیں}}
\end{align*}
 جہاں \عددی{f_1\le f_2} ہے۔ہم اسی طرح کے مفروضے   دیگر محددی سطوح پر \عددی{D}  کے تظلیل قائمہ  کے  لئے فرض کرتے ہیں۔شکل \حوالہء{16.69} دیکھیں۔

اکائی عمودی سمتیہ \عددی{\an=n_1\ai+n_2\aj+n_3\ak} کے اجزاء زاویات \عددی{\alpha}، \عددی{\beta}، اور \عددی{\gamma}   کے کوسائن ہیں جو \عددی{\an} اکائی سمتیات \عددی{\ai}، \عددی{\aj}، اور \عددی{\ak} کے ساتھ بناتا ہے (شکل \حوالہء{16.70})۔  تمام سمتیات اکائی  ہونے کی وجہ سے ایسا ہے۔یوں 
\begin{align*}
n_1&=\an\cdot\ai=|\an||\ai|\cos\alpha=\cos\alpha\\
n_2&=\an\cdot\aj=|\an||\aj|\cos\beta=\cos\beta\\
n_3&=\an\cdot\ak=|\an||\ak|\cos\gamma=\cos\gamma
\end{align*}
لہٰذا
\begin{align*}
\an=(\cos\alpha)\ai+(\cos\beta)\aj+(\cos\gamma)\ak
\end{align*}
اور درج ذیل ہو گا۔
\begin{align*}
\kvec{F}\cdot\an=M\cos\alpha+N\cos\beta+P\cos\gamma
\end{align*}
اجزاء کے روپ میں مسئلہ پھیلاؤ  کہتا ہے درج ذیل  ہو گا۔
\begin{align*}
\iint_S(M\cos\alpha+N\cos\beta+P\cos\gamma)\dif\sigma=\iiint_D\left(\frac{\partial M}{\partial x}+\frac{\partial N}{\partial y}+\frac{\partial P}{\partial z}\right)\dif x\dif y\dif z
\end{align*}
ہم درج ذیل تین  مساویت ثابت کر کے مسئلہ ثابت کرتے ہیں۔
\begin{align}
\iint_S M\cos\alpha\dif \sigma&=\iiint_D \frac{\partial M}{\partial x}\dif x\dif y\dif z\label{مساوات_سمتی_میدان_پھیلا_الف}\\
\iint_S N\cos\beta\dif \sigma&=\iiint_D \frac{\partial N}{\partial y}\dif x\dif y\dif z\label{مساوات_سمتی_میدان_پھیلا_ب}\\
\iint_S P\cos\gamma\dif \sigma&=\iiint_D \frac{\partial P}{\partial z}\dif x\dif y\dif z\label{مساوات_سمتی_میدان_پھیلا_پ}
\end{align}
\موٹا{مساوات \حوالہ{مساوات_سمتی_میدان_پھیلا_پ} کا ثبوت}\quad
ہم بائیں ہاتھ کے سطحی تکمل کو \عددی{xy} مستوی پر \عددی{D} کی قائمہ تظلیل  پر دوہرے تکمل میں تبدیل کر کے مساوات \حوالہ{مساوات_سمتی_میدان_پھیلا_پ} ثابت  کرتے ہیں (شکل \حوالہء{16.71})۔ سطح \عددی{S}   دو سطوح  \عددی{S_2} اور \عددی{S_1} پر مشتمل ہے جہاں \عددی{S_2} بالا حصہ ہے اور اس کی مساوات \عددی{z=f_2(x,y)} ہے اور \عددی{S_1}  نچلا حصہ  ہے اور اس کی مساوات \عددی{z=f_1(x,y)} ہے۔سطح \عددی{S_2} پر باہر رخ عمودی اکائی سمتیہ \عددی{\an}  کا \عددی{\ak} جزو مثبت علامت رکھتا ہے اور
 چونکہ 
\begin{align*}
\dif\sigma=\frac{\dif A}{|\cos\gamma|}=\frac{\dif x\dif y}{\cos\gamma}
\end{align*}
ہے لہٰذا درج ذیل ہو گا (شکل \حوالہء{16.72}  پر نظر ڈالیں)۔
\begin{align*}
\cos\gamma\dif\sigma=\dif x\dif y
\end{align*}
سطح \عددی{S_1} پر \عددی{\an} کا \عددی{\ak} جزو منفی علامت رکھتا ہے لہٰذا  اس سطح پر درج ذیل ہو گا۔
\begin{align*}
\cos\gamma\dif\sigma=-\dif x\dif y
\end{align*}
یوں درج ذیل حاصل ہو گا۔
\begin{align*}
\iint_S P\cos\gamma\dif\sigma&=\iint_{S_2}P\cos\gamma\dif\sigma+\iint_{S_1}P\cos\gamma\dif\sigma\\
&=\iint_{R_{xy}} P(x,y,f_2(x,y))\dif x\dif y-\iint_{R_{xy}}P(x,y,f_1(x,y))\dif x\dif y\\
&=\iint_{R_{xy}}[P(x,y,f_2(x,y))-P(x,y,f_1(x,y))]\dif \dif y\\
&=\iint_{R_{xy}}\left[\int_{f_1(x,y)}^{f_2(x,y)}\frac{\partial P}{\partial z}\dif z\right]\dif x\dif y=\iiint_D\frac{\partial P}{\partial z}\dif z\dif x\dif y
\end{align*}
یوں مساوات \حوالہ{مساوات_سمتی_میدان_پھیلا_پ} ثابت ہوا۔

مساوات \حوالہ{مساوات_سمتی_میدان_پھیلا_الف} اور مساوات \حوالہ{مساوات_سمتی_میدان_پھیلا_ب}  کے ثبوت  اسی طرز پر حاصل ہوں گے؛ یا مساوات \حوالہ{مساوات_سمتی_میدان_پھیلا_پ}  میں بالترتیب \عددی{x}، \عددی{y}، \عددی{z}؛ \عددی{M}، \عددی{N}، \عددی{P}؛ \عددی{\alpha}، \عددی{\beta}،\عددی{\gamma}  ایک دوسرے کی جگہ  رکھ  کر نتائج   لکھیں۔

\جزوحصہء{دیگر خطوں کے لئے مسئلہ پھیلاؤ}
مسئلہ پھیلاؤ کو ایسے خطوں تک وسعت دی جا سکتی ہے جنہیں محدود تعداد  کے مذکورہ  سادہ قسم کے  خطوں  میں تقسیم کیا جا سکتا ہے  یا ایسے خطوں تک  جنہیں  کسی طرح سادہ خطوں کی حد  تصور کیا جا سکتا ہے۔مثال کے طور پر،دو  ہم مرکز کرّات  کے بیچ خطہ   \عددی{D}  ہے  جہاں \عددی{D} کے اندر اور اس کے تحدیدی سطوح  پر \عددی{\kvec{F{}}} کے جزوی تفرق استمراری ہیں۔استوائی مستوی پر \عددی{D} کو دو حصوں میں تقسیم کر کے دونوں نصف حصوں  پر مسئلہ  پھیلاؤ کا اطلاق کریں۔نچلا نصف حصہ  \عددی{D_1} شکل \حوالہء{16.73} میں پیش کیا گیا ہے جس کی تحدیدی سطح \عددی{S_1} ہے جو   بیرونی کرہ، مسطح چھلا، اور   اندرونی کرہ  پر مشتمل ہے۔ مسئلہ پھیلاؤ درج ذیل کہتا ہے۔
\begin{align}
\iint_{S_1}\kvec{F}\cdot\an_1\dif \sigma_1=\iiint_{D_1}\nabla\cdot\kvec{F}\dif H_1\label{مساوات_سمتی_میدان_نچلا_کرہ}
\end{align}
اکائی عمودی سمتیہ \عددی{\an_1}  کا رخ نچلے نصف حصہ \عددی{D_1} کے بیرونی سطح پر مبدا سے دوری جانب ہے، چھلے پر یہ \عددی{\ak} کے برابر ہے، اور اندرونی سطح پر  اس کا رخ مبدا کی جانب ہے۔ اس کے بعد \عددی{D_2} اور اس کی سطح \عددی{S_2}  پر مسئلہ پھیلاؤ کا اطلاق کریں (شکل \حوالہء{16.74})۔
\begin{align}
\iint_{S_2}\kvec{F}\cdot\an_2\dif \sigma_2=\iiint_{D_2}\nabla\cdot\kvec{F}\dif H_2\label{مساوات_سمتی_میدان_بالا_کرہ}
\end{align}
بالائی نصف حصہ \عددی{D_2} کی سطح  \عددی{S_2}   پر چلتے ہوئے  عمادی اکائی سمتیہ \عددی{\an_2} پر نظر رکھیں۔ بیرونی کرہ پر اس کا رخ مبدا سے دور  جانب ہے، چھلا پر یہ \عددی{-\ak} کے برابر ہے، جبکہ اندرونی کرہ پر اس کا رخ مبدا کی جانب ہے۔ مساوات \حوالہ{مساوات_سمتی_میدان_نچلا_کرہ} اور مساوات \حوالہ{مساوات_سمتی_میدان_بالا_کرہ} جمع کرتے ہوئے ، چھلے پر  \عددی{\an_1} اور \عددی{\an_2}   کی علامتیں الٹ ہونے کی بنا ، اس پر تکمل  ایک دوسرے کو کاٹتے ہیں۔یوں درج ذیل نتیجہ حاصل ہو گا
\begin{align*}
\iint_{S}\kvec{F}\cdot\an\dif \sigma=\iiint_{D}\nabla\cdot\kvec{F}\dif H
\end{align*}
جہاں دو کرّات کے بیچ خطہ \عددی{D} ہے، \عددی{D} کی سرحد \عددی{S} دو کرّات پر مشتمل ہے، اور \عددی{S} پر \عددی{D} سے باہر رخ عمودی اکائی سمتیہ \عددی{\an} ہے۔

\ابتدا{مثال}\شناخت{مثال_سمتی_میدان_تکمل_کرہ_بہاؤ}\موٹا{باہر رخ بہاؤ کا حصول}\\
خطہ \عددی{D:\,0< a^2\le x^2+y^2+z^2\le b^2}  کی سرحد کو عبور کرتا ہوا  باہر رخ بہاؤ درج ذیل میدان کے لئے تلاش کریں۔
\begin{align*}
\kvec{F}=\frac{x\ai+y\ai+z\ak}{\rho^3},\quad \rho=\sqrt{x^2+y^2+z^2}
\end{align*}
حل:\quad
خطہ \عددی{D} پر \عددی{\nabla\cdot\kvec{F}} کا تکمل ہمیں بہاؤ دیگا۔درج ذیل لکھتے ہیں
\begin{align*}
\frac{\partial \rho}{\partial x}&=\frac{1}{2}(x^2+y^2+z^2)^{-1/2}(2x)=\frac{x}{\rho}\\
\frac{\partial M}{\partial x}&=\frac{\partial}{\partial x}(x\rho^{-3})=\rho^{-3}-3x\rho^{-4}\frac{\partial \rho}{\partial x}=\frac{1}{\rho^3}-\frac{3x^2}{\rho^5}
\end{align*}
اور اسی طرح درج ذیل بھی ہو گا۔
\begin{align*}
\frac{\partial N}{\partial y}=\frac{1}{\rho^3}-\frac{3y^2}{\rho^5},\quad \frac{\partial P}{\partial z}=\frac{1}{\rho^3}-\frac{3z^2}{\rho^5}
\end{align*}
لہٰذا
\begin{align*}
\kvec{F}\text{پھیلاؤ}=\frac{3}{\rho^3}-\frac{3}{\rho^5}(x^2+y^2+z^2)=\frac{3}{\rho^3}-\frac{3\rho^2}{\rho^5}=0
\end{align*}
اور درج ذیل ہو گا۔
\begin{align*}
\iiint_D \nabla\cdot\kvec{F}\dif H=0
\end{align*}

اس طرح \عددی{D} پر \عددی{\nabla\cdot\kvec{F}} کا تکمل صفر ہے اور \عددی{D} کی سرحد کو  باہر رخ عبور کرتا ہوا  صافی بہاؤ صفر ہو گا۔ اس مثال  سے مزید معلومات بھی حاصل ہوتی ہیں۔ اندرونی کرہ \عددی{S_a}  کو پار کرتا  ہوا جو بہاؤ \عددی{D} سے خارج  ہوتا ہے  کی علامت اس بہاؤ کے الٹ ہو گی جو بیرونی کرہ \عددی{S_b}  کو عبور کرتا ہوا \عددی{D} سے خارج ہوتا ہے (یاد رہے ان بہاؤ کا مجموعہ صفر ہے)۔ یوں مبدا سے دوری کے رخ \عددی{\kvec{F}}   کا بہاؤ جو \عددی{S_a} کو عبور کرتا ہے  \عددی{\kvec{F}}  کے اس بہاؤ کے برابر ہو گا جو   مبدا سے دوری کے رخ \عددی{S_b} کو پار کرتا ہوا \عددی{D} سے خارج ہوتا ہے۔یوں مبدا پر مرکوز کرہ   کو عبور کرتا \عددی{\kvec{F}} کا  بہاؤ کرہ کے رداس پر منحصر نہیں۔ یہ بہاؤ کتنا ہو گا؟

اس کی قیمت جاننے کی خاطر ہم بہاؤ کا تکمل حل کرتے ہیں۔ رداس \عددی{a} کے کرہ  کا باہر رخ اکائی عمودی سمتیہ درج ذیل ہو گا۔
\begin{align*}
\an=\frac{x\ai+y\aj+z\ak}{\sqrt{x^2+y^2+z^2}}=\frac{x\ai+y\aj+z\ak}{a}
\end{align*}
یوں کرہ پر
\begin{align*}
\kvec{F}\cdot\an=\left(\frac{x\ai+y\aj+z\ak}{a^3}\right)\cdot \left(\frac{x\ai+y\aj+z\ak}{a}\right)=\frac{x^2+y^2+z^2}{a^4}=\frac{a^2}{a^4}=\frac{1}{a^2}
\end{align*}
لہٰذا درج ذیل ہو گا۔
\begin{align*}
\iint_{S_a}\kvec{F}\cdot\an\dif\sigma =\frac{1}{a^2}\iint_{S_a}\dif \sigma=\frac{1}{a^2}(4\pi a^2)=4\pi
\end{align*}
مبدا پر مرکوز ہر کرہ سے \عددی{\kvec{F}} کا  باہر  رخ بہاؤ \عددی{4\pi} ہو گا۔
\انتہا{مثال}

\جزوحصہء{گاوس کا قانون: برقناطیسیت کے چار عظیم  قوانین میں سے ایک}
ابھی  مثال \حوالہ{مثال_سمتی_میدان_تکمل_کرہ_بہاؤ} سے بہت کچھ سیکھا جا سکتا ہے۔ برقناطیسی نظریہ میں، مبدا پر موجود  نقطی بار \عددی{q}  درج ذیل برقی میدان پیدا کرتا ہے
\begin{align*}
\kvec{E}(x,y,z)=\frac{1}{4\pi \epsilon_0}\frac{q}{|\kvec{r}|^2}\left(\frac{\kvec{r}}{|\kvec{r}|}\right)=\frac{q}{4\pi \epsilon_0}\frac{\kvec{r}}{|\kvec{r}|^3}=\frac{q}{4\pi\epsilon_0}\frac{x\ai+y\aj+z\ak}{\rho^3}
\end{align*}
جہاں \عددی{\epsilon_0} طبیعی مستقل،  \عددی{\kvec{r}} نقطہ \عددی{(x,y,z)} کا  تعین گر سمتیہ، اور \عددی{\rho=|\kvec{r}|=\sqrt{x^2+y^2+z^2}} ہے۔   مثال \حوالہ{مثال_سمتی_میدان_تکمل_کرہ_بہاؤ}  کی علامتیت  میں  اس کو لکھتے ہیں۔
\begin{align*}
\kvec{E}=\frac{q}{4\pi\epsilon_0}\kvec{F}
\end{align*}

   مثال \حوالہ{مثال_سمتی_میدان_تکمل_کرہ_بہاؤ}  سے  ہم دیکھتے ہیں کہ مبدا پر مرکوز  ہر کرہ سے باہر رخ \عددی{\kvec{E}} کا  بہاؤ \عددی{q/\epsilon_0} ہو گا، لیکن یہ نتیجہ صرف کرہ تک محدود نہیں۔کسی بھی بند سطح \عددی{S} سے، جو مبدا کو گھیرتی ہو (اور جس پر مسئلہ پھیلاؤ کا اطلاق ہوتا ہو)، \عددی{\kvec{E}} کا باہر رخ بہاؤ \عددی{q/\epsilon_0} ہو گا۔اس کی وجہ جاننے کی خاطر مبدا پر مرکوز   اتنے بڑے کرہ \عددی{S_a} کا تصور کریں جو \عددی{S} کو گھیرتا ہو۔جب  \عددی{\rho>0} ہو  درج ذیل ہو گا
\begin{align*}
\nabla\cdot\kvec{E}=\nabla\cdot \frac{q}{4\pi\epsilon_0}\kvec{F}=\frac{q}{4\pi\epsilon_0}\nabla\cdot\kvec{F}=0
\end{align*}
جس کی وجہ سے  \عددی{S} اور \عددی{S_a} کے بیچ خطہ \عددی{D} پر \عددی{\nabla\cdot\kvec{E}}  کا تکمل صفر ہو گا۔لہٰذا مسئلہ پھیلاؤ کے تحت 
\begin{align*}
\iint_{\text{\RL{$D$ کی سرحد}}}\kvec{E}\cdot\an\dif \sigma=0
\end{align*}
ہو گا، اور مبدا سے دوری کے رخ \عددی{S} کو عبور کرتا  ہوا \عددی{\kvec{E}} کا بہاؤ لازماً مبدا سے دوری کے رخ  \عددی{S_a} کو عبور کرتے ہوئے  \عددی{\kvec{E}} کے بہاؤ کے برابر ہو گا، جو \عددی{q/\epsilon_0} ہے۔یہ فقرہ، جو قانون گاوس کہلاتا ہے، یہاں فرض کئے گئے بار سے زیادہ عمومی  تقسیم بار کے لئے بھی درست ہے، جیسا  آپ طبیعیات کے کسی بھی نصابی کتاب میں  پڑھ سکتے ہیں۔
\begin{align*}
\iint_{S}\kvec{E}\cdot\an\dif \sigma&=\frac{q}{\epsilon_0}&&\text{\RL{قانون گاوس}}
\end{align*}

\جزوحصہء{  ماقوا حرکیات  کی  استمراری مساوات}
فرض کریں فضا میں خطہ \عددی{D} بند ، سمت بند سطح \عددی{S} میں محدود  ہے۔ اگر  نقطہ \عددی{(x,y,z)}   اور وقت \عددی{t} پر \عددی{D} میں بہتے ہوئے  سیال کی سمتی رفتاری  میدان  \عددی{\kvec{v}(x,y,z)} ،  اور  کثافت \عددی{\delta=\delta(t,x,y,z)} ہو، اور \عددی{\kvec{F}=\delta\kvec{v}}  ہو، تب ماقوا حرکیات کی  \اصطلاح{استمراری مساوات}\فرہنگ{استمراری مساوات!ماقوا حرکیات}\حاشیہب{continuity equation}\فرہنگ{continuity equation}  درج ذیل کہتی ہے۔
\begin{align*}
\nabla\cdot\kvec{F}+\frac{\partial \delta}{\partial t}=0
\end{align*}
اگر ملوث  تفاعلات کے یک رتبی جزوی تفرقات استمراری ہوں، یہ مساوات  مسئلہ پھیلاؤ سے اخذ  کی جا سکتی ہے۔ آئیں یہ عمل دیکھیں۔

درج ذیل تکمل  ، \عددی{S} کو عبور کر کے \عددی{D} سے  کمیتی اخراج  (اخراج اس لئے کہ \عددی{\an}   باہر رخ عمودی ہے) کی شرح دیتا ہے۔
\begin{align*}
\iint_S \kvec{F}\cdot\an\dif \sigma
\end{align*}
اس کی وجہ جاننے کی خاطر، سطح پر   رقبے کا ٹکڑا \عددی{\Delta \sigma}   لیں (شکل \حوالہء{16.75})۔ قلیل زمانی  وقفہ \عددی{\Delta t}  میں، اس ٹکڑے کو \عددی{\Delta H} حجم عبور کرتا ہے جو تخمیناً اس بیلن کے حجم کے برابر ہو گا جس کا تل \عددی{\Delta \sigma} اور   لمبائی 
\عددی{(\kvec{v}\Delta t)\cdot \an} ہو، جہاں اس ٹکڑے کے کسی ایک نقطے پر سمتی رفتاری سمتیہ \عددی{\kvec{v}} ہے۔
\begin{align*}
\Delta H\approx \kvec{v}\cdot \an\Delta \sigma\Delta t
\end{align*}
اس حجم کی کمیت تقریباً
\begin{align*}
\Delta m\approx \delta \kvec{v}\cdot \an\Delta \sigma\Delta t
\end{align*}
ہو گی، لہٰذا   سے اس ٹکڑے سے گزر کر \عددی{D} سے کمیت کے اخراج کی شرح تخمیناً  درج ذیل ہو گی۔
\begin{align*}
\frac{\Delta m}{\Delta t}\approx \delta \kvec{v}\cdot \an\Delta \sigma
\end{align*}
اس سے درج ذیل تخمین لکھی جا سکتی ہے
\begin{align*}
\frac{\sum \Delta m}{\Delta t}\approx \sum \delta \kvec{v}\cdot \an\Delta \sigma
\end{align*}
جو \عددی{S} کو پار کرتے ہوئے کمیت کی  تخمیناً اوسط شرح ظاہر کرتی ہے۔ آخر میں \عددی{\Delta \sigma\to 0} اور \عددی{\Delta t\to 0} لیتے ہوئے  \عددی{D} سے \عددی{S} کی راہ سے اخراج کی لمحاتی شرح:
\begin{align*}
\frac{\dif  m}{\dif  t}= \iint_S \delta \kvec{v}\cdot \an\dif \sigma
\end{align*}
حاصل ہو گی جو  اس مخصوص سیال  کے لئے درج ذیل لکھی جا سکتی ہے۔
\begin{align*}
\frac{\dif  m}{\dif  t}= \iint_S \kvec{F}\cdot \an\dif  \sigma
\end{align*}

اب فرض کریں سیال میں نقطہ \عددی{Q} پر  مرکوز ٹھوس کرہ \عددی{B} رکھا جاتا ہے۔ کرہ \عددی{B} پر \عددی{\nabla\cdot\kvec{F}} کی اوسط قیمت درج ذیل ہو گی۔
\begin{align*}
\frac{1}{\text{\RL{$B$ کا حجم}}}\iiint_B\nabla\cdot\kvec{F}\dif H
\end{align*}
پھیلاؤ کے استمرار کا ایک پہلو یہ ہے کہ \عددی{B} میں  کسی ایک نقطہ \عددی{P} پر  \عددی{\nabla\cdot\kvec{F}} کی قیمت یہی ہو گی۔یوں درج ذیل ہو گا۔
\begin{align}
(\nabla\cdot\kvec{F})_P&=\frac{1}{\text{\RL{\عددی{B} کا حجم}}}\iiint_B\nabla\cdot\kvec{F}\dif H=\frac{\iint_S\kvec{F}\cdot\an\dif\sigma}{\text{\RL{\عددی{B} کا حجم}}}\notag\\
&=\frac{\text{\RL{سطح \عددی{S} کے ذریعہ \عددی{B} سے کمیتی اخراج کی شرح}}}{\text{\RL{\عددی{B} کا حجم}}}\label{مساوات_سمتی_میدان_کمیتی_شرح_اخراج}
\end{align}
دائیں ہاتھ  کسر اکائی حجم میں کمیت کی کمی بیان کرتا ہے۔

اب مرکز \عددی{Q} کو ایک ہی مقام پر رکھتے ہوئے  \عددی{B} کے رداس کو صفر کے قریب پہنچائیں۔ مساوات  \حوالہ{مساوات_سمتی_میدان_کمیتی_شرح_اخراج} کا بایاں ہاتھ \عددی{(\nabla\cdot\kvec{F})_Q}  کو، جبکہ اس کا دایاں ہاتھ \عددی{(-\partial \delta/\partial t)_Q} کو مرکوز ہو گا۔ان دو حدوں   کی مساویت  ہمیں درج ذیل \اصطلاح{استمراری مساوات }\فرہنگ{استمراری!مساوات}\حاشیہب{continuity equation}\فرہنگ{continuity!equation} دیتی ہے۔
\begin{align*}
\nabla\cdot\kvec{F}=-\frac{\partial \delta}{\partial t}
\end{align*}

استمراری مساوات ہمیں \عددی{\nabla\cdot\kvec{F}} کا مفہوم دیتی ہے: کسی نقطے پر \عددی{\kvec{F}} کا پھیلاؤ اس نقطے پر سیال کی کثافت میں کمی کی شرح دیتا ہے۔

مسئلہ پھیلاؤ:
\begin{align*}
\iint_S\kvec{F}\cdot\an\dif \sigma=\iiint_D\nabla\cdot\kvec{F}\dif H
\end{align*}
اب کہتا ہے کہ خطہ \عددی{D} میں کثافت  کی صافی کمی کی وجہ   کمیت کی  سطح \عددی{S} سے دوسری جانب منتقلی  ہے۔ یوں یہ مسئلہ کمیت کی بقا  بیان کرتا ہے (سوال \حوالہء{31})۔

\جزوحصہء{تکملی مسائل کی یکجا سازی}
